% AUTO-GENERATED by tools/usc_extract.py — do not edit by hand
\begingroup\small
\begin{historicalsource}[{\textbf{Primary text} --- Title~18, \S~981(a) (Civil forfeiture) (OLRC via \texttt{nickvido/us-code}, tag \texttt{annual/2025})}]
\textit{Source:} \url{https://uscode.house.gov/view.xhtml?req=granuleid:USC-prelim-title18&num=0&edition=prelim}\par\medskip\textbf{(a)}

(1) The following property is subject to forfeiture to the United States:

  (A) Any property, real or personal, involved in a transaction or attempted transaction in violation of section 1956, 1957 or 1960 of this title, or any property traceable to such property.

  (B) Any property, real or personal, within the jurisdiction of the United States, constituting, derived from, or traceable to, any proceeds obtained directly or indirectly from an offense against a foreign nation, or any property used to facilitate such an offense, if the offense—

    (i) involves trafficking in nuclear, chemical, biological, or radiological weapons technology or material, or the manufacture, importation, sale, or distribution of a controlled substance (as that term is defined for purposes of the Controlled Substances Act), or any other conduct described in section 1956(c)(7)(B);

    (ii) would be punishable within the jurisdiction of the foreign nation by death or imprisonment for a term exceeding 1 year; and

    (iii) would be punishable under the laws of the United States by imprisonment for a term exceeding 1 year, if the act or activity constituting the offense had occurred within the jurisdiction of the United States.

  (C) Any property, real or personal, which constitutes or is derived from proceeds traceable to a violation of section 215, 471, 472, 473, 474, 476, 477, 478, 479, 480, 481, 485, 486, 487, 488, 501, 502, 510, 542, 545, 656, 657, 670, 842, 844, 1005, 1006, 1007, 1014, 1028, 1029, 1030, 1032, or 1344 of this title or any offense constituting “specified unlawful activity” (as defined in \href{https://uscode.house.gov/view.xhtml?req=granuleid:USC-prelim-title18-section1956/c/7&num=0&edition=prelim}{section 1956(c)(7) of this title}), or a conspiracy to commit such offense.

  (D) Any property, real or personal, which represents or is traceable to the gross receipts obtained, directly or indirectly, from a violation of—

    (i) section 666(a)(1) (relating to Federal program fraud);

    (ii) section 1001 (relating to fraud and false statements);

    (iii) section 1031 (relating to major fraud against the United States);

    (iv) section 1032 (relating to concealment of assets from conservator or receiver of insured financial institution);

    (v) section 1341 (relating to mail fraud); or

    (vi) section 1343 (relating to wire fraud),
    if such violation relates to the sale of assets acquired or held by the the 11 So in original. Federal Deposit Insurance Corporation, as conservator or receiver for a financial institution, or any other conservator for a financial institution appointed by the Office of the Comptroller of the Currency or the National Credit Union Administration, as conservator or liquidating agent for a financial institution.

  (E) With respect to an offense listed in subsection (a)(1)(D) committed for the purpose of executing or attempting to execute any scheme or artifice to defraud, or for obtaining money or property by means of false or fraudulent statements, pretenses, representations or promises, the gross receipts of such an offense shall include all property, real or personal, tangible or intangible, which thereby is obtained, directly or indirectly.

  (F) Any property, real or personal, which represents or is traceable to the gross proceeds obtained, directly or indirectly, from a violation of—

    (i) section 511 (altering or removing motor vehicle identification numbers);

    (ii) section 553 (importing or exporting stolen motor vehicles);

    (iii) section 2119 (armed robbery of automobiles);

    (iv) section 2312 (transporting stolen motor vehicles in interstate commerce); or

    (v) section 2313 (possessing or selling a stolen motor vehicle that has moved in interstate commerce).

  (G) All assets, foreign or domestic—

    (i) of any individual, entity, or organization engaged in planning or perpetrating any any 1 Federal crime of terrorism (as defined in section 2332b(g)(5)) against the United States, citizens or residents of the United States, or their property, and all assets, foreign or domestic, affording any person a source of influence over any such entity or organization;

    (ii) acquired or maintained by any person with the intent and for the purpose of supporting, planning, conducting, or concealing any Federal crime of terrorism (as defined in section 2332b(g)(5) 22 So in original. A second closing parenthesis probably should appear. against the United States, citizens or residents of the United States, or their property;

    (iii) derived from, involved in, or used or intended to be used to commit any Federal crime of terrorism (as defined in section 2332b(g)(5)) against the United States, citizens or residents of the United States, or their property; or

    (iv) of any individual, entity, or organization engaged in planning or perpetrating any act of international terrorism (as defined in section 2331) against any international organization (as defined in section 209 of the State Department Basic Authorities Act of 1956 (\href{https://uscode.house.gov/view.xhtml?req=granuleid:USC-prelim-title22-section4309/b&num=0&edition=prelim}{22 U.S.C. 4309(b)}) or against any foreign Government.33 So in original. Probably should not be capitalized. Where the property sought for forfeiture is located beyond the territorial boundaries of the United States, an act in furtherance of such planning or perpetration must have occurred within the jurisdiction of the United States.

  (H) Any property, real or personal, involved in a violation or attempted violation, or which constitutes or is derived from proceeds traceable to a violation, of section 2339C of this title.

  (I) Any property, real or personal, that is involved in a violation or attempted violation, or which constitutes or is derived from proceeds traceable to a prohibition imposed pursuant to section 104(a) of the North Korea Sanctions and Policy Enhancement Act of 2016.

(2) For purposes of paragraph (1), the term “proceeds” is defined as follows:

  (A) In cases involving illegal goods, illegal services, unlawful activities, and telemarketing and health care fraud schemes, the term “proceeds” means property of any kind obtained directly or indirectly, as the result of the commission of the offense giving rise to forfeiture, and any property traceable thereto, and is not limited to the net gain or profit realized from the offense.

  (B) In cases involving lawful goods or lawful services that are sold or provided in an illegal manner, the term “proceeds” means the amount of money acquired through the illegal transactions resulting in the forfeiture, less the direct costs incurred in providing the goods or services. The claimant shall have the burden of proof with respect to the issue of direct costs. The direct costs shall not include any part of the overhead expenses of the entity providing the goods or services, or any part of the income taxes paid by the entity.

  (C) In cases involving fraud in the process of obtaining a loan or extension of credit, the court shall allow the claimant a deduction from the forfeiture to the extent that the loan was repaid, or the debt was satisfied, without any financial loss to the victim.
\end{historicalsource}
\endgroup
